\documentclass{article}
\usepackage{graphicx} % Required for inserting images
\usepackage{lipsum}   % For filler text
\usepackage[hidelinks]{hyperref}
\usepackage{amsmath, amssymb}
\usepackage{xcolor}

% For hyperlinks and navigation (custom environment placeholder)
\newenvironment{callout}
{

    % begin code
}
{
    % end code
}

\title{Additive Manufacturing}
\author{Julia Kim}
\date{November 2025}

\begin{document}

\maketitle

\section{Additive Manufacturing}

Additive Manufacturing (AM) is a process that builds parts layer-by-layer directly from a digital model, rather than subtracting from a bulk material like conventional manufacturing processes. Additive manufacturing enables complex geometry, lightweighting, and part consolidation. 3D Printing is one category of additive manufacturing.


\section{3D Printing Process}

\begin{itemize}
    \item 1. 3D CAD model is converted to STL (Standard Tessellation Language)
    \item 2. Model is sliced into 2D layers
    \item 3. Each layer is fabricated sequentially
    \item 4. Support material may be required, depending on geometry and process.
\end{itemize}

\begin{callout}
Standard Tessellation Language (STL) represents surfaces using triangular facets
Tessellation: surface represented by triangles
image
\end{callout}

\section{Advantages and Disadvantages}
There are a wide range of 3D Printing technologies that each have their own unique properties. Between them, you can optimize certain qualities with other tradeoffs, but these refer to 3D Printing as a whole.

Advantages
\begin{itemize}
    \item Easy and quick to change part design and produce a new version (rapid prototyping)
    \item Limited skills or training needed to make parts
    \item Reduced individual manufacturing time with no prep time required
    \item Cost of production unique parts in small quantities is relatively low
    \item Some parts can be made pre-assembled
\end{itemize}
Disadvantages
\begin{itemize}
    \item Part tolerances/surface finish:
        \begin{itemize}
            \item "Staircase" Effect on sloped surfaces from slicing and layer thickness
            \item Shrinkage or distortion of manufactured parts from uneven cooling
            \item Holes come out smaller and posts bigger than CAD dimensions
        \end{itemize}
    \item Limited variety of possible materials
    \item Post-processing of parts may be required
    \item Cost and time does not improve with increased volume
    \item Part size is generally limited
    \item Porosity and mechanical property challenges
    \item Material anisotropy, i.e. mechanical properties within layers vary from those between layers
\end{itemize}

\section{Types of 3D Printing Processes}

\begin{enumerate}
    \item Material Extrusion (MEX)
    \begin{itemize}
        \item Fused Deposition Modeling ($(FDM)^{TM}$)
        \item Fused Filament Fabrication (FFF)
    \end{itemize}
    \item Vat Polymerization (VAT)
    \begin{itemize}
        \item Stereolithography (SLA)
        \item Digital Light Processing (DLP)
        \item Continuous Digital Light Processing (CLIP)
        \item Liquid Crystal Display (LCD)
    \end{itemize}
    \item Material Jetting
    \begin{itemize}
        \item Poly Jet
        \item Nanoparticle Jetting
        \item Drop on Demand (DOD)
    \end{itemize}
    \item Directed Energy Deposition (DED)
    \begin{itemize}
        \item Laser Engineered Net Shaping ($LENS^{TM}$)
        \item Electron Beam Additive Manufacturing (EBAM)
        \item Wire Arc Additive Manufacturing (WAAM)
        \item Friction Stir Additive Manufacturing (FSAM)
    \end{itemize}
    \item Powder Bed Fusion (PBF)
    \begin{itemize}
        \item Selective Laser Sintering (SLS)
        \item Selective Laser Melting (SLM)
        \item Direct Metal Laser Sintering (DMLS)
        \item Electron Beam Melting (EBM)
        \item Selective Heat Sintering (SHS)
        \item High-Speed Sintering
        \item Multi-Jet Fusion (MJF)
    \end{itemize}
    \item Binder Jetting
    \item Sheet Lamination 
    \begin{itemize}
        \item Laminate Object Manufacturing (LOM)
        \item Selective Lamination Composite Object Manufacturing
        \item Plastic Sheet Lamination
        \item Selective Deposition Lamination
        \item Composite Based Additive Manufacturing
        \item Ultrasonic Consolidation
        \item Ultrasonic Additive Manufacturing (UAM)
    \end{itemize}
\end{enumerate}

\subsection{Material Extrusion (MEX)}
This is what most people picture when they think of 3D Printing
\textbf{Also known as:} Fused Deposition Modeling (FDM), Fused Filament Fabrication (FFF)

\begin{itemize}
    \item Heated thermoplastic filament extruded through a nozzle, layers fusing together in cooler air
    \item Machine deposits layers by moving nozzle horizontally, tracing 2D cross sections, then moving vertically and repeating
    \begin{itemize}
        \item 3+ DOF machines exist for non-planar "slices", but much more complex and less common
    \end{itemize}
    \item Support structure needed for steeply overhanging parts
    \item Main Materials: PLA, ABS, PC, Nylon
    \item Typical layer thickness: $0.1$–$0.3$ mm
    \item Surface finish is generally not as good as other methods
\end{itemize}

\subsubsection{MEX Numerical Example}
Given:
\begin{itemize}
    \item Nozzle diameter $d = 0.4$ mm
    \item Layer height $h = 0.25$ mm
    \item Extrusion speed $v = 50$ mm/s
    \item Delay between layers = 5 s
    \item Part: $30 \times 30 \times 30$ mm cube
\end{itemize}

Volumetric flow rate:
\[
Q = v \cdot d \cdot h
= 50 \times 0.4 \times 0.25
= \boxed{5\ \text{mm}^3/\text{s}}
\]

Number of layers:
\[
N = \frac{30}{0.25} = 120
\]

Total inactive time:
\[
t_{\text{idle}} = 120 \times 5 = 600 \text{ s} = 10 \text{ min}
\]

\paragraph{(a) Solid Cube}

Volume:
\[
V = 30^3 = 27{,}000 \ \text{mm}^3
\]

Time to extrude:
\[
t = \frac{V}{Q} = \frac{27{,}000}{5} = 5400 \text{ s} = 90 \text{ min}
\]

Total production time:
\[
t_{\text{total}} = 90 + 10 = \boxed{100 \text{ min}}
\]

\paragraph{(b) 25\% Infill}

Interior dimensions:
\[
30 - 2(0.4) = 29.2 \text{ mm}, \quad \text{height} \approx 29.5 \text{ mm}
\]

Interior volume:
\[
V_{\text{int}} = 29.2^2 \times 29.5
\]

Material used:
\[
V_{\text{total}} = (30^3) - (29.2^2 \cdot 29.5)\cdot 0.75 = 8{,}135 \ \text{mm}^3
\]

Extrusion time:
\[
t = \frac{8{,}135}{5} = 1627.3 \text{ s} \approx 27.1 \text{ min}
\]

Total production time:
\[
t_{\text{total}} = 27.1 + 10 = \boxed{37.1 \text{ min}}
\]

\paragraph{Traverse Path Method (Check)}

Tracks per layer:
\[
\frac{30}{0.4} = 75
\]

Travel distance per layer:
\[
75 \times 30 = 2250 \ \text{mm}
\]

Time per layer:
\[
\frac{2250}{50} = 45 \text{ s}
\]

Total extrusion time:
\[
45 \times 120 = 5400 \text{ s} = 90 \text{ min}
\]

Total production time:
\[
90 + 10 = 100 \text{ min}
\]

\subsection{Vat Polymerization (VAT)}

\begin{itemize}
    \item UltraViolet laser cures UV-curable polymer (or "resin") dissolved in a solvent- Liquid Deposition Process
    \item Laser may be positioned above or below resin tank (top-down vs bottom-up)
    \item Support structures required when overhang angle exceeds $45^\circ$, since cured resin is denser than uncured
    \item SLA, DLP, CLIP are common variants
    \item Typical layer thickness: $0.025$–$0.1$ mm
\end{itemize}

\subsubsection{VAT Build Time Model}

Time for layer $i$:
\[
T_i = \frac{A_i}{vD} + t_d
\]

Total build time:
\[
T = \sum_{i=1}^{n} T_i
\]

\subsubsection{VAT Numerical Example}
A rectangular part with dimensions $20 \text{ mm} \times 10 \text{ mm} \times 10 \text{ mm}$ is fabricated using a vat photopolymerization process.

Given:
\begin{itemize}
    \item Layer thickness $t = 0.2$ mm
    \item Part height $H = 10$ mm
    \item Laser scanning speed $v = 1000$ mm/s
    \item Laser spot diameter $D = 0.2$ mm
    \item Delay per layer (stage movement in $z$-direction) $t_d = 10$ s
    \item Cross-sectional area per layer $A = 200$ mm$^2$
\end{itemize}

\textbf{Find:}
\begin{enumerate}
    \item Number of layers
    \item Time to write one layer (including delay)
    \item Total fabrication time
    \item Volumetric method approximation for fabrication time
\end{enumerate}

\textbf{Solution:}

\begin{enumerate}
    \item Number of layers:
\[
n = \frac{H}{t} = \frac{10}{0.2} = 50
\]
    \item Time to write one layer:
\[
T_i = \frac{A}{v \cdot D} + t_d
= \frac{200}{1000 \times 0.2} + 10
= 11 \ \text{s}
\]
    \item Total fabrication time:
\[
T = n \cdot T_i = 50 \times 11 = \boxed{550 \ \text{s}}
\]
    \item Volumetric method:

Total part volume:
\[
V_{\text{part}} = A \cdot H = 200 \times 10 = 2000 \ \text{mm}^3
\]

Volumetric build rate:
\[
\dot{V} = D \cdot t \cdot v = (0.2)(0.2)(1000)
\]

Total time:
\[
T = \frac{V_{\text{part}}}{D \cdot t \cdot v} + n \cdot t_d
= \frac{2000}{(0.2)(0.2)(1000)} + 50 \times 10
= 50 + 500
= \boxed{550 \ \text{s}}
\]
\end{enumerate}

\subsection{Material Jetting}

\begin{itemize}
    \item Individual droplets of photopolymer deposited layer-by-layer and UV-cured
    \item Machines typically have 2 sets of print heads, one for main material and one for support, with multiple nozzles on each head
    \item Uses viscous thermosetting resin that is cured by high intensity UV, with a secondary support material that is water-soluble.
    \item Multiple materials and colors possible
    \item Typical layer thickness ~0.02 mm
\end{itemize}

\subsection{Powder Bed Fusion (PBF)}

\begin{itemize}
    \item Laser or electron beam selectively melts heat-fusible powder (thermoplastics, metals, or ceramics)
    \item After each layer, a new layer of loose powder is spread across the surface
    \item Unfused powder acts as support structure for overhanging parts, but powder in a 10 mm thick shell around the part is thermally damaged
    \item Undamaged powder can be poured out of complete part to be reused
    \item Typical layer thickness: $50$–$100\ \mu$m
\end{itemize}

\subsection{Directed Energy Deposition (DED)}

\begin{itemize}
    \item The material (usually a metal powder) and the energy for fusion (usually laser heating) are simultaneously focused at the same location
    \item The metal powder is projected onto a surface where the laser directly melts the powder in place
    \item Can be combined with in-situ machining, combining both machining and deposition in a single machine
    \item Suitable for repair and large structures
    \item Tool change or part shifting required, requires specialized toolpath planning
\end{itemize}

\subsection{Binder Jetting}

\begin{itemize}
    \item Liquid adhesive is selectively deposited to join powder materials (gypsum or starch), which are deposited in thin layers
    \item Parts are weak before post-processing
\end{itemize}

\subsection{Laminated Object Manufacturing}

\begin{itemize}
    \item The input is a roll of thin sheet material (paper, plastic, or metal) coated with an adhesive
    \item Successive layers are cut by a laser in the required shape and stacked to build up the part
\end{itemize}

\subsection{Design Guidance for AM}

\begin{itemize}
    \item 1. Minimize unnecessary material
    \item 2. Choose a print orientation that is aware of the direction of anisotropy, mechanical properties, surface finish, roundness of holes, support materials, etc.
    \item 3. Reduce support structures
    \item 4. Iterate steps 1-3
\end{itemize}

\end{document}
